% !TeX root = ../../Book2.tex
\section{调和函数}
\label{b2:sec:1301}

热传导达到平衡而区域内部没有热源时，温度应当满足 Laplace 方程。这样的函数在一个小球中心的值由球面或球内平均决定，局部不能凭空出现新的峰值。这种平均结构还会产生光滑性、比较估计和唯一性，远比单独观察一个二阶微分等式更强。

\Needspace{6\baselineskip}
\begin{definition}\label{b2:def:13-harmonic}
设整数 \(d\geq1\)，\(\Omega\subset\mathbb R^d\) 为开集，\(\mathbb K\in\{\mathbb R,\mathbb C\}\)。若 \(u\in C^2(\Omega;\mathbb K)\) 且 \(\Delta u=0\)，则称 \(u\) 为 \(\Omega\) 上的\term[diaohe-hanshu]{调和函数}[harmonic function]。除另作说明外，本章先取实值函数；复值函数的调和性按实部和虚部分别理解。
\end{definition}

一次函数调和，二次函数 \(x_1^2-x_2^2\) 也调和，但 \(|x|^2\) 的 Laplace 等于 \(2d\)，并不调和。一维调和函数恰为仿射函数；本章涉及球面核和基本解的主要计算固定 \(d\geq2\)。沿用 \(\omega_d=|B(0,1)|\)、\(\sigma_{d-1}=|\partial B(0,1)|=d\omega_d\)，分别表示体积与表面积。

\subsection{平均值性质}
\label{b2:subsec:130101}

\Needspace{6\baselineskip}
\begin{theorem}[调和函数的平均值性质]\label{b2:thm:13-mean-value}
设整数 \(d\geq1\)，\(\Omega\subset\mathbb R^d\) 开、\(a\in\Omega\)、\(R>0\)。若实值或复值函数 \(u\) 在 \(\Omega\) 上调和，且 \(\overline{B(a,R)}\subset\Omega\)，则对 \(0<r\leq R\)，有
\[
 u(a)=\dfrac1{\sigma_{d-1}r^{d-1}}\int_{\partial B(a,r)}u\,\mathrm{d}S
 =\dfrac1{\omega_dr^d}\int_{B(a,r)}u\,\mathrm{d}x.
\]
\end{theorem}
\begin{proof}
将球面平均写到固定单位球面上，记
\[
 M(r)=\dfrac1{\sigma_{d-1}}\int_{\partial B(0,1)}u(a+r\theta)\,\mathrm{d}S_\theta.
\]
对固定的紧半径区间可在积分下求导。换回半径为 \(r\) 的球面，再用散度定理，得到
\[
 M'(r)=\dfrac1{\sigma_{d-1}r^{d-1}}
          \int_{\partial B(a,r)}\partial_\nu u\,\mathrm{d}S
       =\dfrac1{\sigma_{d-1}r^{d-1}}\int_{B(a,r)}\Delta u\,\mathrm{d}x=0.
\]
所以 \(M\) 为常数，而 \(r\downarrow0\) 时由连续性有 \(M(r)\to u(a)\)，得到球面公式。最后用极坐标沿 \(0<s<r\) 积分，球内积分等于 \(\int_0^r\sigma_{d-1}s^{d-1}u(a)\,\mathrm{d}s=\omega_dr^du(a)\)。
\end{proof}

这个\term[pingjunzhi-xingzhi]{平均值性质}[mean-value property]不含导数，因而也可以反过来用于识别调和性。

\Needspace{6\baselineskip}
\begin{proposition}[平均值性质的逆命题]\label{b2:prop:13-converse-mean}
设整数 \(d\geq1\)，\(\Omega\subset\mathbb R^d\) 为开集，\(\mathbb K\in\{\mathbb R,\mathbb C\}\)，\(u\in C(\Omega;\mathbb K)\)。若每个闭球包含于 \(\Omega\) 时，球面平均均等于中心值，则 \(u\) 光滑且调和。若 \(u\) 连续并在分布意义下满足 \(\Delta u=0\)，也有同样结论。
\end{proposition}
\begin{proof}
取非负径向磨光函数 \(\rho\in C_c^\infty(B(0,1))\)，积分为 \(1\)，并令 \(\rho_\varepsilon(x)=\varepsilon^{-d}\rho(x/\varepsilon)\)。球面平均公式按径向权积分，给出 \(u*\rho_\varepsilon=u\)，只要核的支集不触及边界。对任一闭包包含于 \(\Omega\) 的球，可取同一个小 \(\varepsilon\) 使这一等式在该球附近成立。左侧光滑，故 \(u\) 局部光滑。

现在对 \(u(a+r\theta)\) 作二阶 Taylor 展开。球面上 \(\theta_j\) 的平均为零，\(\theta_i\theta_j\) 在 \(i\ne j\) 时平均为零，\(\theta_j^2\) 的平均为 \(1/d\)。因此球面平均为
\[
 u(a)+\dfrac{r^2}{2d}\Delta u(a)+o(r^2).
\]
它恒等于 \(u(a)\)，除以 \(r^2\) 再令 \(r\downarrow0\) 可得 \(\Delta u(a)=0\)。

若起初只知分布方程，内部磨光 \(u_\varepsilon=u*\rho_\varepsilon\) 满足经典方程 \(\Delta u_\varepsilon=0\)，因为分布求导与磨光交换。连续性保证 \(u_\varepsilon\to u\) 在紧集上一致成立。对任意固定内部闭球，令 \(\varepsilon\downarrow0\)，便将 \(u_\varepsilon\) 的球面平均公式传给 \(u\)，再应用前半部分。
\end{proof}

球面平均与球内平均的归一化不可混用。若将球面的未归一化面积积分直接视为平均，会额外带入 \(\sigma_{d-1}r^{d-1}\)。同样，平均值性质针对整个小球而非某个坐标方向；只要求每条坐标线上的中点值等于端点平均，得到的是另一种更强且不同的条件。

有时只知道每个中心附近的平均性质，允许使用的半径还随中心变化。这比刚才对全部内部球的假设弱，不能直接为整个邻域选定同一个卷积尺度。

\Needspace{6\baselineskip}
\begin{proposition}[局部平均值性质的逆命题]\label{b2:prop:13-local-converse-mean}
设整数 \(d\geq1\)，\(\Omega\subset\mathbb R^d\) 为开集，\(\mathbb K\in\{\mathbb R,\mathbb C\}\)，\(u\in C(\Omega;\mathbb K)\)。若对每个 \(a\in\Omega\)，存在 \(r_a>0\)，使 \(\overline{B(a,r_a)}\subset\Omega\)，并且对所有 \(0<r<r_a\) 有
\[
 u(a)=\frac1{\sigma_{d-1}r^{d-1}}
           \int_{\partial B(a,r)}u\,\mathrm{d}S,
 \qquad \sigma_{d-1}=|S^{d-1}|,
\]
则 \(u\) 在 \(\Omega\) 内光滑且调和。
\end{proposition}

证明延后至\subsecref{b2:subsec:13-ball-boundary}，在那里先建立球上连续 Dirichlet 问题的解，再用它将局部均值性质推广到全部内部球。本节此前已经证明的\cref{b2:prop:13-converse-mean}不需要这一步构造。

\subsection{最大值原理}
\label{b2:subsec:130102}

\Needspace{6\baselineskip}
\begin{theorem}[调和函数的强最大值原理]\label{b2:thm:13-strong-maximum}
设整数 \(d\geq1\)，\(\Omega\subset\mathbb R^d\) 为连通开集，\(u\) 为 \(\Omega\) 上的实值调和函数。若 \(u\) 在 \(\Omega\) 内达到它在该区域上的最大值或最小值，则 \(u\) 为常数。
\end{theorem}
\begin{proof}
设最大值 \(M\) 在 \(a\) 处达到。对一个内部小球，非负连续函数 \(M-u\) 的球内平均为零，所以它在整个球上为零：若某点取正值，连续性会给出一个正测度区域，其积分为正。故集合 \(\{u=M\}\) 在 \(\Omega\) 中开；由连续性它也闭，且非空。连通性使它等于 \(\Omega\)。最小值情形对 \(-u\) 使用同一论证。
\end{proof}

这项\term[diaohe-qiang-zuidazhi-yuanli]{调和函数强最大值原理}[strong maximum principle for harmonic functions]中的“达到”不能改成只有有限上确界。例如 \(u(x_1,x_2)=e^{x_1}\cos x_2\) 在半平面 \(x_1<0\) 上调和且有界，上确界为 \(1\)，却不在内部达到。连通性也不可省去，否则不同连通分支可以取不同常数。

\Needspace{6\baselineskip}
\begin{corollary}[边界最大值与比较]\label{b2:cor:13-boundary-maximum}
设整数 \(d\geq1\)，\(\Omega\subset\mathbb R^d\) 为非空有界开集，实值 \(u\in C(\overline\Omega)\) 在 \(\Omega\) 内调和，则
\[
 \min_{\partial\Omega}u\leq u(x)\leq\max_{\partial\Omega}u.
\]
若两个这样的调和函数在边界满足 \(u\leq v\)，则区域内也满足 \(u\leq v\)。
\end{corollary}
\begin{proof}
连续函数在紧集 \(\overline\Omega\) 上达到最大值。若它严格超过边界最大值，只能在某个内部连通分支达到；强最大值原理使函数在该分支恒定，沿分支内点趋于其边界便与边界条件矛盾。对 \(-u\) 得到下界，对 \(u-v\) 得到比较结论。这里不需要边界光滑。
\end{proof}

复值调和函数也有模的最大值原理。若 \(|u|\leq M\) 且 \(|u(a)|=M\)，选单位复数 \(c\) 使 \(cu(a)=M\)。实值调和函数 \(\operatorname{Re}(cu)\leq M\) 在内部达到最大值，故恒为 \(M\)；再由 \(|cu|\leq M\) 得虚部为零。于是 \(u\) 为常数。实函数的大小比较和复函数的模比较应分别使用，不能把复数直接放进大小不等式。

\subsection{Harnack 不等式}
\label{b2:subsec:130103}

最大值原理比较内部和边界。若函数非负，则平均值性质还可以比较两个内部点，而无需事先知道边界数据。

\Needspace{6\baselineskip}
\begin{lemma}\label{b2:lem:13-local-harnack}
设整数 \(d\geq1\)，\(a\in\mathbb R^d\)、\(R>0\)，实值 \(u\ge0\) 在 \(B(a,R)\) 上调和。当 \(|x-a|\leq R/4\) 时，
\[
 2^{-d}u(a)\leq u(x)\leq(3/2)^d u(a).
\]
\end{lemma}
\begin{proof}
包含关系 \(B(a,R/4)\subset B(x,R/2)\subset B(a,3R/4)\) 与非负性给出两个积分比较。对三只球分别使用平均值性质，再除以 \(|B(x,R/2)|\)，即得体积比 \(2^{-d}\) 与 \((3/2)^d\)。所有闭球都留在调和域内部，端点情形可由稍缩半径后取极限得到。
\end{proof}

\Needspace{6\baselineskip}
\begin{theorem}[Harnack]\label{b2:thm:13-harnack}
设整数 \(d\geq1\)，\(a\in\mathbb R^d\)、\(R>0\)。若实值 \(u\ge0\) 在 \(B(a,R)\) 上调和，\(0<r<R\)，则存在只依赖 \(d,r/R\) 的常数 \(C\)，使
\[
 \sup_{B(a,r)}u\leq C\inf_{B(a,r)}u.
\]
更一般地，若 \(\Omega\subset\mathbb R^d\) 为连通开集、实值 \(u\ge0\) 在 \(\Omega\) 上调和、\(K\Subset\Omega\)，则存在只依赖 \(K,\Omega\) 的相同类型比较常数。
\end{theorem}
\begin{proof}
连接内球中两点的线段完全留在内球。令 \(\delta=R-r\)，将线段分成长度至多 \(\delta/4\) 的有限段；每个分点为中心的半径 \(\delta\) 开球都在原域内。逐段使用局部引理，一步的双向比较均可用 \(2^d\) 控制，段数至多 \(\lceil8r/\delta\rceil+1\)。连乘即得只依赖半径比的常数，最后对两点分别取上、下确界。

一般连通开集中的两点可以用域内折线连接：从一个固定点可用折线到达的集合在域中既开又闭，因此是全域。每条有限折线与边界保持正距离，可以被有限条上述小球链覆盖。对紧集 \(K\)，再选有限个终端小球覆盖它，合并有限条连接链，便得到对全部 \(K\) 统一的常数。比较只用于非负函数，故每次积分包含关系都保留方向。
\end{proof}

这称为\term[harnack-budengshi]{Harnack 不等式}[Harnack's inequality]。非负条件是实质性的：函数 \(u(x)=x_1\) 调和，但一个包含原点的球中同时有正、负值，不满足这种正比例比较。对非负调和函数，如果在一点取零，Harnack 或强最小值原理都会使它在整个连通域为零。

\Needspace{6\baselineskip}
\begin{corollary}[Harnack 收敛原理]\label{b2:cor:13-harnack-convergence}
设整数 \(d\geq1\)，\(\Omega\subset\mathbb R^d\) 为连通开集，\((u_n)\) 为 \(\Omega\) 上逐点递增的实值调和函数列，\(a\in\Omega\)。若 \(u_n(a)\) 在某一点有有限上界，则 \(u_n\) 在每个紧集上一致收敛到调和函数；若该点的值趋于 \(+\infty\)，则在每个紧集上一致趋于 \(+\infty\)。
\end{corollary}
\begin{proof}
第一种情形中，对 \(m>n\) 将 Harnack 应用于非负函数 \(u_m-u_n\)。它在紧集上的上确界由常数乘 \(u_m(a)-u_n(a)\) 控制，故为一致 Cauchy 列。极限连续，球面平均公式按局部一致极限保留，由已证明的\cref{b2:prop:13-converse-mean}知极限调和。第二种情形对 \(u_n-u_1\) 作反向的紧集比较，其下确界随 \(u_n(a)-u_1(a)\) 趋于无穷。
\end{proof}

\subsection{内部正则性与 Liouville 定理}
\label{b2:subsec:130104}

\Needspace{6\baselineskip}
\begin{proposition}[内部导数估计]\label{b2:prop:13-interior-derivatives}
设整数 \(d\geq1\)，\(a\in\mathbb R^d\)、\(R>0\)，实值或复值 \(u\) 在 \(B(a,2R)\) 上调和，则对每个多重指标 \(\alpha\)，
\[
 \sup_{B(a,R)}|\partial^\alpha u|
 \leq C_{d,\alpha}R^{-d-|\alpha|}\int_{B(a,2R)}|u|\,\mathrm{d}x.
\]
\end{proposition}
\begin{proof}
先将球面平均按径向磨光权积分。对 \(x\in B(a,R)\)，有 \(u(x)=\int\rho_R(x-y)u(y)\,\mathrm{d}y\)，核支集限于 \(B(a,2R)\)。将导数移到光滑核，\(\|\partial^\alpha\rho_R\|_\infty=R^{-d-|\alpha|}\|\partial^\alpha\rho\|_\infty\)，得到所求估计。等式本身也再次证明了调和函数的任意阶光滑性。
\end{proof}

\Needspace{6\baselineskip}
\begin{proposition}[局部可积调和函数的代表]\label{b2:prop:13-harmonic-l1-representative}
设整数 \(d\geq1\)，\(\Omega\subset\mathbb R^d\) 为开集，\(\mathbb K\in\{\mathbb R,\mathbb C\}\)。若 \(h\in L^1_{\mathrm{loc}}(\Omega;\mathbb K)\)，且 \(\Delta h=0\) 在分布意义下成立，则 \(h\) 有唯一的光滑调和代表。因而内部导数估计也适用于这个代表，右端仍使用原等价类的绝对积分。
\end{proposition}
\begin{proof}
固定 \(\overline{B(a,3R)}\subset\Omega\)，用半径 \(\varepsilon<R/2\) 的磨光核构造 \(h_\varepsilon\)。分布求导与卷积交换，所以 \(h_\varepsilon\) 在 \(B(a,2R)\) 上光滑且调和。局部的 \(L^1\) 逼近恒等给出 \(h_\varepsilon\to h\) 于 \(L^1(B(a,2R))\)。对两个磨光尺度之差使用内部估计，得到
\[
 \sup_{B(a,R)}|\partial^\alpha(h_\varepsilon-h_\delta)|
 \leq C_{d,\alpha}R^{-d-|\alpha|}
       \|h_\varepsilon-h_\delta\|_{L^1(B(a,2R))}\longrightarrow0.
\]
因此任意有限阶导数都在内部小球上一致 Cauchy；逐次使用微积分基本定理，极限是光滑函数，其各阶导数正是相应的一致极限。二阶导数收敛使极限仍调和，零阶 \(L^1\) 收敛又将它识别为 \(h\) 的几乎处处代表。

这些内部球覆盖 \(\Omega\)。两个局部代表在重叠处几乎处处相同，而连续函数若在一个开集上几乎处处为零，便处处为零。因此局部代表能够拼合，且光滑代表唯一。最后对任一内部球使用前一命题，代表与 \(h\) 的绝对积分相同，故估计保留。
\end{proof}

这一步把分布方程与经典调和性连接起来，所用条件只是局部可积性。磨光函数的每个导数虽各自存在，仅有 \(L^1\) 收敛却原本不足以传递导数；真正提供一致控制的是调和函数的内部估计。

\Needspace{6\baselineskip}
\begin{theorem}[Liouville]\label{b2:thm:13-liouville}
设整数 \(d\geq1\)，\(\mathbb K\in\{\mathbb R,\mathbb C\}\)，\(u\) 为整个 \(\mathbb R^d\) 上的 \(\mathbb K\) 值调和函数。若 \(u\) 全局有界，则 \(u\) 为常数。实值 \(u\) 只要有全局上界或下界，也必为常数。若存在 \(C>0\) 和非负整数 \(m\)，使全部 \(x\in\mathbb R^d\) 满足 \(|u(x)|\leq C(1+|x|)^m\)，则 \(u\) 为次数不超过 \(m\) 的调和多项式。
\end{theorem}
\begin{proof}
有界情形对任意点使用一阶内部估计，让 \(R\to\infty\)，得到梯度处处为零。若实函数上有界，取 \(v=M-u\geq0\)。内部估计中的绝对积分等于 \(v\) 的积分，而平均值性质给出 \(\int_{B(a,2R)}v=\omega_d(2R)^dv(a)\)。因此 \(|\nabla v(a)|\leq C_dR^{-1}v(a)\to0\)。下有界时改用 \(u-M\)。

在多项式增长条件下，固定 \(a\)，取 \(|\alpha|=m+1\)。内部估计的右侧为 \(O(R^{-1})\)，故全部 \(m+1\) 阶偏导数为零。对连接原点与任意点的线段使用一元 Taylor 公式，余项中的 \(m+1\) 阶方向导数为零，所以 \(u\) 恰等于在原点的 \(m\) 阶 Taylor 多项式。该多项式仍满足原来的调和方程。
\end{proof}

上述结论统称调和函数的\term[diaohe-liouville]{Liouville 定理}[Liouville's theorem for harmonic functions]。增长条件不能删除，例如 \(x_1\) 是非常数整调和函数；仅有界于某个真子域时也不能套用全空间结论。后面处理半空间唯一性时，必须同时利用边界条件和适当的增长类。
